\section{Choosing the next move}\label{appE:sec:next-move}
Mathematics can look like a long list of tricks to memorize, but it helps far more to recognize the kind of situation you are in. Four situations recur in this book, and each suggests a natural next move.

\subsection*{When a general claim looks doubtful}
Look for a counterexample before you try to prove the claim. Start with small and boundary cases: zero, the empty set, the smallest allowed input, two objects that happen to be equal. A useful counterexample is rarely found by trying cases at random; it is built to break the step that the claim seems to need.
\begin{itemize}
\item Proposition~\ref{thm:compconverse} shows that $g\circ f$ can be injective although $g$ is not. There $f$ maps $\{0,1\}$ into $\{0,1,2\}$ and never reaches $2$, while $g$ sends $2$ to the same output as $1$. The composition applies $g$ only to outputs of $f$, so it never meets this collision.
\item In Section~\ref{ch18:sec:guess}, the guess that small update errors still give convergence to the fixed point fails for a constant bias: the map $x\mapsto qx$ on $\RR$, with $0<q<1$, has the fixed point $0$, but adding the same $\eta>0$ at every update moves the limit to $\eta/(1-q)$, however small $\eta$ is.
\end{itemize}
Each counterexample also showed what a correct statement must take into account: the outputs that $f$ actually reaches, and errors that never die out.

\subsection*{When you need to show that something exists}
When a direct construction is out of reach, ask what a failure would force. The book turns that question into a proof in three ways.
\begin{itemize}
\item \emph{Search and repair.} The search of Chapter~\ref{ch:matching} either finds an augmenting path, which enlarges the matching, or fails and produces a group of students who together accept too few tasks (Proposition~\ref{ch08:prop:shortage}), so that no matching gives every student a task. Either outcome tells us something, which is the point of Habit~7 (Section~\ref{ch08:sec:shortage}).
\item \emph{An extremal choice.} Choose an object that is as large or as small as possible, and show that it already has the property you want. A longest directed path in a nonempty finite acyclic oriented graph ends at a sink, because an arrow leaving its last vertex would give either a longer path or a directed cycle (Theorem~\ref{thm:sink}). Similarly, if a statement about natural numbers fails somewhere, choose the smallest number at which it fails: the statement holds at every smaller number, and this often leads to a contradiction (Section~\ref{ch05:sec:not-circular}).
\item \emph{An average.} A random variable $X$ on a finite probability space satisfies $X(\omega)\ge\E X$ at some outcome $\omega$ of positive probability, since outcomes of probability zero contribute nothing to $\E X$ and a weighted average cannot exceed all the values it averages (Proposition~\ref{ch13:prop:above-average}). Theorem~\ref{thm:cut} uses this: labeling the vertices of a graph with $m$ edges at random gives a cut of expected size $m/2$, so some labeling gives a cut of size at least $m/2$, although the proof never says which. Habit~8 (Section~\ref{ch13:sec:markov}) describes this approach.
\end{itemize}
To choose among them, ask: Can a partial solution be repaired one step at a time? Is there a natural quantity to make as large or as small as possible? Is there an average you can compute exactly?

\subsection*{When the full answer is out of reach}
Ask what you really need to know. Often it is much less than the full answer.
\begin{itemize}
\item To know how close a point $z$ is to the fixed point of a contraction, you need not find the fixed point: the residual estimate of Theorem~\ref{thm:banach} bounds the error by $d(z,Tz)/(1-q)$, and $Tz$ is usually easy to compute. In Section~\ref{ch12:sec:example}, this certified that $x_2=10/27$ lies within $19/729$ of the fixed point of $T(x)=(x^2+1)/3$ on $[0,1]$ before the fixed point had been found.
\item To show that no circulant Hadamard matrix of order $12$ exists, you do not need the complete answer claimed by the research of Section~\ref{ch14:sec:research}. The necessary condition of Proposition~\ref{thm:circulantnecessary} is enough: an order greater than $1$ must have the form $4u^2$ for a positive integer $u$, and $12$ lies strictly between the first two such numbers, $4$ and $16$.
\end{itemize}
The risk in settling for less is answering an easier question than the one you were asked. If you replace the question by an easier one, say so, and compare the new version with the contract of the original, as Habit~2 (Section~\ref{sec:magnitude}) asks. Proposition~\ref{thm:circulantnecessary}, for instance, settles order $12$, but as an answer to ``which orders are possible?'' it is only partial: besides $1$ and $4$, which do occur, it still allows $16,36,64,\ldots$, and says nothing about whether any of them occurs.

\subsection*{When a bound appears}
Ask whether the constant in the bound could be lowered and, if not, which example prevents any improvement. Such an example usually attains the bound, or comes as close to it as we like, so also ask when the bound is attained (Section~\ref{ch09:sec:taste}).
\begin{itemize}
\item The Cauchy--Schwarz inequality $|u\cdot v|\le\norm{u}_2\norm{v}_2$ holds with equality exactly when $v=0$ or $u$ is a multiple of $v$ (Section~\ref{ch10:sec:cs}). Because equality occurs for nonzero vectors, such as $u=v\ne0$, the constant factor $1$ on the right cannot be made smaller.
\item In the constant-bias example of Section~\ref{ch18:sec:sharp}, the error bound of Theorem~\ref{thm:inexact}, part~2, holds with equality at every step: with bias $\eta>0$, the errors are exactly $\eta(1-q^n)/(1-q)$, which approach $\eta/(1-q)$ as $n$ grows. So the factor $1/(1-q)$ is sharp. Any smaller factor would give a bound below the true error once $n$ is large enough, so no cleverer proof can lower it.
\end{itemize}
